%----------
%	CONCLUSIONES
%----------

Este trabajo ha abordado la detección automática de lenguaje sexista en la prensa
española escrita, a partir del corpus y la metodología del proyecto IRIS, con el fin de asistir el análisis
experto de la representación mediática de las mujeres. La propuesta central ha sido
empaquetar la metodología experta como \textit{Agent Skills} cargables bajo demanda,
sobre una arquitectura de agentes especializadas, una por cada variable del libro de
códigos (V25, V26, V30, V33 y V35). Para valorar su aportación real, esa arquitectura
(nivel B1) se comparó con una condición de control que recibía la misma metodología
inyectada en el \textit{prompt} (nivel B0), sobre un corpus de 1.313 noticias anotadas
y cuatro modelos de lenguaje, tres accesibles por API y uno local.

La evaluación experimental, interpretada en detalle en la Sección~\ref{sec:discussion},
deja varias conclusiones que aquí se recogen de forma sintética. La primera es que
empaquetar la metodología como \textit{Agent Skills} \textbf{no mejora el acuerdo por
sí mismo}, ya que solo un modelo registra una ganancia de entidad, la de
\texttt{gemini} en \texttt{masc\_generico}, donde su $\kappa$ sube 0,188 puntos,
mientras en los demás los movimientos se quedan en unas pocas milésimas o van hacia
abajo. La arquitectura amplifica, por tanto, la capacidad que cada modelo ya trae en
lugar de conferirla, y la herramienta de recuperación en vivo (RAG) no aporta una
ganancia que justifique su coste. La segunda es de alcance metodológico: en una tarea con
clases muy desequilibradas, la exactitud y $\kappa$ pueden moverse en sentidos opuestos,
por lo que reportar solo la exactitud llevaría a la conclusión contraria y se impone
informar de exactitud, $\kappa$ y prevalencia de forma conjunta. La tercera es que el
bajo acuerdo absoluto responde a \textbf{dos causas que conviene no confundir}. Una es
un techo de la propia tarea, documentado en \texttt{sexismo\_discurso}, que dos equipos
marcan en el 11,5\,\% y en el 65,2\,\% de las piezas y que dentro de un mismo equipo va
del 28,3\,\% al 97,8\,\%, reflejo de lo difícil que resulta percibir el fenómeno incluso
a quien lo busca. La otra es una limitación de los modelos: en \texttt{lenguaje\_sexista} los dos
equipos marcan la variable en la mayoría de sus piezas, el 74,6\,\% y el 85,5\,\%
respectivamente, y aun así tres de los cuatro modelos apenas la detectan, con un
recall prácticamente nulo. El peso de cada una cambia según la variable, y reparte responsabilidades, de
manera que el desacuerdo recoge a la vez la dificultad de percibir el fenómeno y el
error de detección del sistema. Cabe
precisar, además, que los modelos quedan, salvo \texttt{gemini} en
\texttt{masc\_generico}, por debajo de la persona que menos marca en términos de \textit{recall}, lo que responde a
una capacidad de detección menor y no a un criterio distinto. La cuarta es que un modelo local
de pesos abiertos, \texttt{gemma4:e4b}, compite variable a variable con los servicios
comerciales accesibles por API, a los que supera en \texttt{asimetria\_mujer\_hombre}, con la cautela de que esa
ventaja descansa en un solo acierto, y a dos de los cuales adelanta en
\texttt{masc\_generico}, preservando la privacidad del material y sin
tarifas de inferencia, a cambio de trasladar el coste al tiempo de cómputo. Realizar esa comparación sobre la misma tarea y el mismo corpus, con un flujo de procesamiento agnóstico a la
empresa proveedora, ha sido uno de los ejes metodológicos del trabajo.

En cuanto al cumplimiento de los objetivos de la Sección~\ref{sec:objectives_motivation},
el balance es el siguiente. El objetivo general (OG1) se ha satisfecho al investigar y
evaluar de forma sistemática una arquitectura basada en agentes y LLM para detectar
variables cuantitativas y cualitativas de representación mediática, delimitando con evidencia empírica
tanto sus prestaciones como sus límites. El diseño conceptual del sistema, con agentes
especializadas y sus flujos de análisis y verificación (OE1), y las estrategias de
\textit{prompting} y herramientas asociadas, criterios, plantillas y contextos expertos
empaquetados como \textit{skills} (OE2), se desarrollaron en el
Capítulo~\ref{implementation} y se validaron en el Capítulo~\ref{results_and_discussion}.
Las técnicas de representación y recuperación de información, incluida la RAG y la
similitud léxica (OE3), se implementaron y evaluaron, con el hallazgo de que la recuperación
en vivo no produce una ganancia consistente que justifique su coste. La comparación
entre modelos y su idoneidad para detectar sesgos explícitos e implícitos (OE4)
constituye el núcleo del Capítulo~\ref{results_and_discussion}. Por último, la propuesta
de directrices para la integración editorial y la supervisión humana (OE5) se concreta
en las directrices que se enuncian al cierre de este capítulo y en las líneas futuras
que se detallan después.

De todo ello se desprende, en coherencia con la discusión, que el sistema encaja mejor
como \textbf{asistente que como decisor}. Su salida tiene dos planos, y solo uno se ha
podido evaluar. El primero es la etiqueta, el plano cuantitativo, que se contrasta con
la anotación de referencia y es el que se mide aquí. El segundo es el marcado de
evidencias y la explicación del veredicto, es decir, el plano cualitativo, esa segunda dimensión
que un análisis de contenido experto también necesita. Este último no cuenta todavía
con un marcado experto de referencia con el que compararse, de modo que su validación
queda como línea futura. El valor del sistema es, en definitiva, práctico y subordinado
al criterio experto: ordenar el corpus, señalar los casos que merecen una segunda
lectura y aliviar el trabajo repetitivo de quienes codifican.

El alcance de estas conclusiones está acotado por varias limitaciones. El corpus
procede de la prensa española y se analizan cinco de las quince variables de lenguaje
del libro de códigos, de modo que la generalización a otros medios, idiomas o variables
queda por comprobar. La evaluación se ha realizado en régimen \textit{zero-shot} y
sobre modelos ligeros y económicos, sin explorar el ajuste fino (\textit{fine-tuning})
ni modelos de mayor capacidad. La recuperación sobre las guías se implementó con similitud léxica (TF-IDF),
y una variante semántica podría comportarse de otro modo. Dos de las variables
descansan además sobre muy pocos casos positivos, 81 en
\texttt{asimetria\_mujer\_hombre} y 131 en \texttt{denominacion\_sexualizada}, de modo
que sus cifras se juegan en un puñado de piezas y admiten menos confianza que las
demás. Tampoco ha podido evaluarse el salto semántico, ya que la anotación de
referencia no emplea ese valor en ninguna pieza del corpus. Como cada pieza fue
codificada por una sola persona, no fue posible calcular la fiabilidad
intercodificadora, por lo que el techo de la tarea se acota de forma indirecta y con incertidumbre significativa.
Por último, y como ya se ha señalado, solo se ha evaluado el plano cuantitativo de la
salida. Estas limitaciones no invalidan los hallazgos, pero delimitan su alcance y
orientan las líneas de trabajo que se proponen a continuación.

De cara a una futura integración del sistema en entornos editoriales reales (OE5), el
trabajo permite formular siete directrices:
\begin{itemize}
    \item Desplegarlo como asistente que ordena y prioriza el corpus, y no como decisor
    automático.
    \item Preservar la trazabilidad de cada predicción, con sus evidencias y explicación,
    para que la analista pueda auditarla, aceptarla o corregirla.
    \item Mantener a la persona en el bucle, de modo que su criterio prevalezca sobre la
    salida del modelo, en especial en las variables en las que el sistema deja sin
    señalar la mayor parte de los positivos.
    \item Decidir su uso variable a variable, asignando a cada una el modelo que mejor se
    comporta en ella, en lugar de desplegar una única configuración para las cinco.
    \item Orientar cualquier ajuste futuro a elevar el \textit{recall}, aun a costa de la
    precisión, porque en postedición un falso negativo es un caso que nadie revisará y
    un falso positivo cuesta unos segundos de descarte.
    \item Favorecer el despliegue local cuando el material sea sensible, por razones de
    privacidad y de coste.
    \item Informar del desempeño con métricas conscientes del desequilibrio, esto es, acuerdo
    y prevalencia además de exactitud, para no sobrestimar la fiabilidad del sistema.
\end{itemize}
